\documentclass[12pt,a4paper]{article}

% ── Packages ──
\usepackage[utf8]{inputenc}
\usepackage[T1]{fontenc}
\usepackage{mathpazo}
\usepackage[margin=2.5cm]{geometry}
\usepackage{graphicx}
\usepackage{booktabs}
\usepackage{longtable}
\usepackage{array}
\usepackage{multirow}
\usepackage{amsmath,amssymb}
\usepackage{float}
\usepackage{enumitem}
\usepackage{xcolor}
\usepackage{hyperref}
\usepackage{cleveref}
\usepackage{caption}
\usepackage{subcaption}
\usepackage{listings}
\usepackage{algorithm}
\usepackage{algpseudocode}
\usepackage{natbib}
\usepackage{setspace}
\usepackage{parskip}
\usepackage{fancyhdr}
\usepackage{titlesec}

% ── Configuration ──
\hypersetup{
    colorlinks=true,
    linkcolor=blue!60!black,
    citecolor=blue!60!black,
    urlcolor=blue!60!black
}

\definecolor{codebg}{rgb}{0.95,0.95,0.97}
\definecolor{codeframe}{rgb}{0.8,0.8,0.85}
\lstset{
    backgroundcolor=\color{codebg},
    frame=single,
    rulecolor=\color{codeframe},
    basicstyle=\ttfamily\small,
    breaklines=true,
    numbers=left,
    numberstyle=\tiny\color{gray},
    keywordstyle=\color{blue!70!black}\bfseries,
    commentstyle=\color{green!50!black}\itshape,
    stringstyle=\color{red!60!black},
    showstringspaces=false,
    tabsize=4
}

\pagestyle{fancy}
\fancyhf{}
\fancyhead[L]{\small BC3409 Individual Assignment}
\fancyhead[R]{\small Earnings Call Analyzer}
\fancyfoot[C]{\thepage}
\renewcommand{\headrulewidth}{0.4pt}

\onehalfspacing

% ── Title ──
\title{%
    \vspace{-1cm}
    \textbf{Earnings Call Analyzer} \\[0.3cm]
    \Large A Fully Local NLP and RAG System for \\
    Earnings Call Transcript Analysis \\[0.5cm]
    \large BC3409 --- AI in Accounting and Finance \\
    Individual Assignment
}
\author{Seunghwan Choi}
\date{March 2026}

\begin{document}

\maketitle
\thispagestyle{empty}

\begin{abstract}
This report presents the Earnings Call Analyzer, a fully local natural language processing (NLP) and Retrieval-Augmented Generation (RAG) system that transforms raw earnings call transcripts into actionable financial intelligence. The system ingests PDF or text transcripts, automatically identifies speakers and their roles, and applies a multi-model sentiment analysis pipeline comprising Loughran--McDonald financial dictionaries, VADER sentiment intensity analysis, and FinBERT transformer-based classification. A two-layer risk detection framework combines keyword-based extraction (255 terms across 10 financial risk categories with negation handling and severity scoring) with semantic risk detection using sentence-transformer embeddings and differential scoring against risk and safe anchor sentences. The RAG subsystem employs hybrid retrieval (FAISS dense search + BM25 sparse search) fused via Reciprocal Rank Fusion and re-ranked by a cross-encoder, with answers generated by a locally-hosted Llama~3.2 language model. The system is trained on 13,603 earnings events spanning 1,776 companies and evaluated across retrieval quality (Recall@3 = 86.9\%, MRR = 100\%), answer quality (citation rate = 83.3\%), and financial predictive power (portfolio sort spread = 0.72\% per trade). An interactive Streamlit dashboard provides multi-quarter analysis with full explainability. The entire pipeline runs locally with zero API costs, ensuring complete data privacy.
\end{abstract}

\newpage
\tableofcontents
\newpage

% ══════════════════════════════════════════════════════════════
\section{Introduction}
\label{sec:intro}
% ══════════════════════════════════════════════════════════════

\subsection{Motivation}

Earnings calls are quarterly events where corporate executives present financial results and field questions from sell-side analysts. A typical transcript spans 60--80 pages of dialogue containing management tone, forward guidance, risk disclosures, and competitive commentary---qualitative signals that complement quantitative financial data. However, the volume and complexity of these transcripts make systematic analysis prohibitively time-consuming for human analysts.

Prior academic research has established that textual features extracted from earnings calls carry predictive information for post-earnings stock returns \citep{loughran2011liability, price2012earnings, chen2018sentimental}. Critically, \citet{price2012earnings} demonstrate that the Q\&A section of earnings calls is more informative than prepared remarks, as management responses to live analyst questioning reveal information not available in scripted presentations. \citet{loughran2011liability} show that finance-specific lexicons significantly outperform general-purpose sentiment dictionaries in financial text analysis.

However, the relationship between text-derived sentiment and stock returns is subtle. Raw sentiment scores are typically statistically insignificant for predicting returns \citep{chen2018sentimental}. What matters is the \emph{change} in sentiment quarter-over-quarter, the divergence between management and analyst tone, and whether sentiment deteriorates when executives face live questions. These are the specific signals this system is designed to extract.

\subsection{Objectives}

This project aims to build a comprehensive, end-to-end system that:
\begin{enumerate}[noitemsep]
    \item Automatically parses earnings call transcripts, identifying speakers, roles, and sections.
    \item Applies three complementary sentiment models with research-backed signal extraction.
    \item Detects financial risks using both keyword-based and semantic approaches with full explainability.
    \item Provides a RAG-based Q\&A interface for natural language querying of transcript content.
    \item Links NLP features to post-earnings stock returns through event studies and ML models.
    \item Runs entirely locally, requiring no API keys or external data transmission.
\end{enumerate}

\subsection{Scope and Data}

The system is trained and evaluated on a dataset comprising:
\begin{itemize}[noitemsep]
    \item \textbf{1,776 unique companies} across all market sectors
    \item \textbf{13,603 earnings events} spanning 2017Q4--2024Q4
    \item \textbf{30,235 text chunks} embedded in the FAISS retrieval index
    \item \textbf{39 engineered NLP features} per earnings event for ML modelling
\end{itemize}

The data originates from two sources: a Kaggle earnings call dataset providing structured transcripts for 49 S\&P~500 companies, and Motley Fool transcripts expanding coverage to 1,776 companies. Market data (stock prices and S\&P~500 benchmark) is sourced from Yahoo Finance via the \texttt{yfinance} library.


% ══════════════════════════════════════════════════════════════
\section{System Architecture}
\label{sec:architecture}
% ══════════════════════════════════════════════════════════════

The system follows a modular, five-layer architecture designed for extensibility and local execution:

\begin{enumerate}[noitemsep]
    \item \textbf{Ingestion Layer}: PDF/text parsing, speaker detection, role classification, section segmentation, and sentence-aware chunking.
    \item \textbf{NLP Processing Layer}: Three-model sentiment analysis, two-layer risk detection, readability metrics, Q\&A dynamics features, and forward-looking statement extraction.
    \item \textbf{Retrieval Layer}: FAISS dense indexing, BM25 sparse indexing, Reciprocal Rank Fusion, and cross-encoder re-ranking.
    \item \textbf{Generation Layer}: Llama~3.2 (3B) via Ollama for RAG-based question answering with citation tracking.
    \item \textbf{Presentation Layer}: Multi-page Streamlit dashboard with interactive visualisations.
\end{enumerate}

\Cref{tab:tech-stack} summarises the key technologies and models used at each layer.

\begin{table}[H]
\centering
\caption{Technology stack and model specifications.}
\label{tab:tech-stack}
\small
\begin{tabular}{@{}llll@{}}
\toprule
\textbf{Component} & \textbf{Model / Library} & \textbf{Specification} \\
\midrule
Sentence Embeddings & \texttt{all-MiniLM-L6-v2} & 384-dim, 22M params \\
Sentiment (Transformer) & \texttt{ProsusAI/finbert} & BERT-base, 110M params \\
Sentiment (Dictionary) & Loughran--McDonald & Finance-specific lexicon \\
Sentiment (Rule-based) & VADER & Valence-aware, compound score \\
Cross-Encoder & \texttt{ms-marco-MiniLM-L-6-v2} & Re-ranking model \\
LLM & \texttt{llama3.2:3b} via Ollama & 3B params, local inference \\
Dense Index & FAISS \texttt{IndexFlatIP} & Exact cosine similarity \\
Sparse Index & BM25Okapi & \texttt{rank\_bm25} library \\
Dashboard & Streamlit & Multi-page application \\
Market Data & \texttt{yfinance} & Yahoo Finance API \\
\bottomrule
\end{tabular}
\end{table}


% ══════════════════════════════════════════════════════════════
\section{Data Pipeline}
\label{sec:data-pipeline}
% ══════════════════════════════════════════════════════════════

\subsection{Data Sources}

Transcripts are sourced from two complementary datasets:

\begin{enumerate}[noitemsep]
    \item \textbf{Kaggle Earnings Call Dataset}: Structured transcripts for 49 S\&P~500 companies across 28 quarters (2018Q1--2024Q4), stored as a Parquet file with columns for ticker, quarter, date, and transcript text.
    \item \textbf{Motley Fool Transcripts}: 18,755 transcripts across 2,876 tickers in plain-text format, organised as per-company directories with per-quarter \texttt{.txt} files.
\end{enumerate}

A unified manifest (\texttt{data/manifest.json}) indexes all 1,776 unique tickers with standardised metadata. Ticker normalisation handles historical changes (e.g., \texttt{FB} $\to$ \texttt{META}, \texttt{GOOG} $\to$ \texttt{GOOGL}).

\subsection{Speaker Detection}
\label{sec:speaker-detection}

Identifying individual speakers in unstructured transcript text is a non-trivial parsing challenge, as formatting conventions vary significantly across transcript providers. The system employs a cascade of 11 regular expression patterns ordered by decreasing specificity:

\begin{enumerate}[noitemsep]
    \item \texttt{Name -- Title} (double dash with title)
    \item \texttt{Name --- Title} (em dash with title)
    \item \texttt{Name (Title)} at end of line
    \item \texttt{Name, Title:} at end of line
    \item \texttt{[Name]} bracket format
    \item All-caps names (\texttt{JOHN SMITH})
    \item Simple \texttt{Name:} format
    \item Various inline variants
\end{enumerate}

A blacklist filter removes false positives from section headers (e.g., ``Prepared Remarks'', ``Question-and-Answer Session'', ``Forward-Looking Statements''). Speaker names are constrained to $\leq 40$ characters and must contain $\geq 2$ words or match known role keywords.

\subsection{Role Classification}
\label{sec:role-classification}

Each detected speaker is classified into one of four role categories: \textbf{Management} (CEO, CFO, COO, CTO, VP), \textbf{Analyst}, \textbf{Investor Relations (IR)}, or \textbf{Operator}. Classification proceeds through a four-tier cascade:

\begin{enumerate}[noitemsep]
    \item \textbf{Title-based}: Extract role from inline title strings using a dictionary of role keywords (e.g., ``Chief Executive Officer'' $\to$ CEO).
    \item \textbf{Context scanning}: Search the first 20\% of the transcript for role mentions within an 80-character window of speaker names.
    \item \textbf{Heuristic inference}: IR-specific cues (``investor relations'', ``on the call with me'', ``joining us today'') and analyst cues (investment bank names: Goldman Sachs, Morgan Stanley, Citi, etc.).
    \item \textbf{ML fallback}: A TF-IDF + Logistic Regression classifier (\texttt{RoleClassifier}) trained on labelled chunks for speakers remaining unclassified.
\end{enumerate}

The \texttt{RoleClassifier} uses a TF-IDF vectoriser with \texttt{max\_features}=5000, \texttt{ngram\_range}=(1,2), \texttt{min\_df}=2, \texttt{max\_df}=0.95, and \texttt{sublinear\_tf}=True, feeding into a Logistic Regression model with \texttt{class\_weight}=``balanced'' and multinomial cross-entropy loss.

\subsection{Section Detection}

Transcripts are segmented into \textbf{Prepared Remarks} and \textbf{Q\&A} sections. Section boundaries are identified using 40+ detection patterns including:
\begin{itemize}[noitemsep]
    \item Prepared remarks cues: ``prepared remarks'', ``opening remarks'', ``welcome to'', ``good morning''
    \item Q\&A cues: ``questions and answers'', ``q\&a session'', ``open the line for questions''
\end{itemize}

A constraint prevents Q\&A detection before the first 20\% of the transcript to avoid false triggers from operator scripts.

\subsection{Chunking Strategy}

Text is chunked using a sentence-aware approach with the following parameters:
\begin{itemize}[noitemsep]
    \item \textbf{Chunk size}: 400 tokens
    \item \textbf{Chunk overlap}: 50 tokens (last sentences from previous chunk)
    \item \textbf{Minimum chunk size}: 50 tokens
\end{itemize}

Chunks are split at sentence boundaries (detected by \texttt{[.!?]} patterns) to preserve semantic coherence. Chunks are merged within speaker turns and split at sentence boundaries if exceeding 2,000 characters. Boilerplate text (safe harbour disclaimers, forward-looking statement warnings, copyright notices) is removed via regex patterns prior to chunking.

The output is a structured DataFrame with columns: \texttt{chunk\_id}, \texttt{ticker}, \texttt{quarter}, \texttt{speaker}, \texttt{role}, \texttt{section}, \texttt{text}, and \texttt{chunk\_index}, stored as Parquet at \texttt{data/processed/all\_chunks.parquet}.


% ══════════════════════════════════════════════════════════════
\section{Sentiment Analysis}
\label{sec:sentiment}
% ══════════════════════════════════════════════════════════════

Sentiment analysis employs three complementary models, each capturing different linguistic dimensions. The rationale for this multi-model approach is that no single method reliably captures the full spectrum of sentiment in financial text \citep{loughran2011liability, hutto2014vader}.

\subsection{Loughran--McDonald Financial Dictionary}

The Loughran--McDonald (LM) lexicon \citep{loughran2011liability} is a finance-specific word list that addresses the well-documented failure of general-purpose sentiment dictionaries in financial contexts. Words like ``liability'', ``tax'', and ``cost'' carry negative connotations in general English but are neutral in financial discourse. The LM dictionary provides curated word lists for six sentiment categories:

\begin{itemize}[noitemsep]
    \item \textbf{Positive}: Words indicating favourable outcomes
    \item \textbf{Negative}: Words indicating adverse outcomes (e.g., ``impairment'', ``restructuring'', ``write-down'')
    \item \textbf{Uncertainty}: Hedging language (e.g., ``approximately'', ``contingent'', ``uncertain'')
    \item \textbf{Litigious}: Legal terminology
    \item \textbf{Constraining}: Words indicating limitations
\end{itemize}

Scores are computed as token-level proportions:
\begin{equation}
    \text{LM Net Score} = \frac{\text{positive\_count} - \text{negative\_count}}{\text{total\_words}}
    \label{eq:lm-score}
\end{equation}
\begin{equation}
    \text{LM Uncertainty Score} = \frac{\text{uncertainty\_count}}{\text{total\_words}}
    \label{eq:lm-uncertainty}
\end{equation}

\subsection{VADER Sentiment Intensity Analyzer}

VADER (Valence Aware Dictionary and sEntiment Reasoner) \citep{hutto2014vader} is a rule-based model optimised for social media and informal text. It handles punctuation emphasis (e.g., ``great!!!'' scores higher than ``great''), capitalisation, degree modifiers, and contrastive conjunctions. VADER produces a compound score in $[-1, 1]$ alongside positive, negative, and neutral proportions.

While not finance-specific, VADER complements the LM dictionary by capturing \emph{how} something is said (emphasis, intensity) rather than \emph{what} is said (financial terminology).

\subsection{FinBERT Transformer Model}

FinBERT \citep{araci2019finbert} is a BERT-base model (\texttt{ProsusAI/finbert}, 110M parameters) fine-tuned on financial news and communications. Unlike dictionary methods, FinBERT captures contextual semantics---it understands that ``not bad'' is positive and that ``declining to comment'' differs from ``declining revenue.''

Configuration:
\begin{itemize}[noitemsep]
    \item \textbf{Max sequence length}: 512 tokens
    \item \textbf{Batch size}: 8 (configurable)
    \item \textbf{Truncation}: Enabled, with chunks $>$2,000 characters truncated prior to tokenisation
    \item \textbf{Output}: Three-class probabilities (positive, negative, neutral) and the predicted label
\end{itemize}

Sentence-level FinBERT scoring is performed by splitting each chunk at sentence boundaries and scoring individually, then averaging back to the chunk level.

\subsection{Aggregation and Signal Extraction}

Beyond raw sentiment scores, the system computes research-backed derivative signals:

\begin{enumerate}[noitemsep]
    \item \textbf{Management--Analyst Divergence}: $\text{mgmt\_sentiment} - \text{analyst\_sentiment}$. When management tone is significantly more positive than analyst tone, this may indicate management over-optimism or credibility concerns.
    \item \textbf{Prepared--Q\&A Divergence}: $\text{prepared\_sentiment} - \text{qa\_sentiment}$. A large drop from prepared remarks to Q\&A suggests that management's scripted optimism does not withstand analyst scrutiny \citep{price2012earnings}.
    \item \textbf{Quarter-over-Quarter Momentum}: $\text{sentiment}_{t} - \text{sentiment}_{t-1}$, computed per ticker. Sentiment momentum is a stronger predictor of returns than sentiment levels \citep{chen2018sentimental}.
\end{enumerate}


% ══════════════════════════════════════════════════════════════
\section{Risk Detection}
\label{sec:risk}
% ══════════════════════════════════════════════════════════════

Risk detection employs a two-layer hybrid architecture combining precision-oriented keyword matching with recall-oriented semantic detection. This design ensures that both explicit risk mentions (``supply chain disruption'') and implicit risk signals (``our pipeline is weaker than expected'') are captured.

\subsection{Risk Taxonomy}

The system defines 10 financial risk categories, codified in \texttt{configs/risk\_taxonomy.json}:

\begin{table}[H]
\centering
\caption{Risk taxonomy with keyword counts and example terms.}
\label{tab:risk-taxonomy}
\small
\begin{tabular}{@{}lrl@{}}
\toprule
\textbf{Category} & \textbf{Keywords} & \textbf{Example Terms} \\
\midrule
Demand Risk & 28 & declining demand, order cancellations \\
Margin Risk & 26 & margin pressure, cost inflation \\
Supply Chain Risk & 24 & supply shortage, logistics challenges \\
Regulatory Risk & 28 & litigation, antitrust, GDPR \\
Competition Risk & 24 & market share loss, commoditization \\
FX / Macro Risk & 27 & strong dollar, recession risk \\
Liquidity Risk & 24 & cash burn, debt maturity \\
Tech Execution Risk & 24 & product delay, cybersecurity breach \\
Labour Risk & 24 & talent shortage, layoffs \\
Geopolitical Risk & 26 & trade war, tariffs, sanctions \\
\midrule
\textbf{Total} & \textbf{255} & \\
\bottomrule
\end{tabular}
\end{table}

\subsection{Layer 1: Keyword-Based Detection}

The keyword layer performs exact substring matching of 255 financial risk terms against each text chunk. Three contextual enhancements improve precision:

\subsubsection{Negation Handling}

A negation window of 3 words preceding each keyword match checks for 19 negation terms (e.g., ``no'', ``not'', ``never'', ``without'', ``mitigated'', ``addressed'', ``overcame'', ``offset''). Negated detections are flagged with \texttt{severity="negated"} and excluded from risk counts. This prevents false positives from statements like ``no supply chain issues'' or ``we mitigated the margin pressure.''

\subsubsection{Severity Scoring}

A 10-word context window (5 before, 5 after the keyword) is scanned for intensity modifiers:
\begin{itemize}[noitemsep]
    \item \textbf{Intensifiers} (20 terms): ``significant'', ``major'', ``severe'', ``critical'', ``substantial'', ``unprecedented'', ``devastating'', etc.
    \item \textbf{Diminishers} (20 terms): ``minor'', ``slight'', ``marginal'', ``modest'', ``limited'', ``manageable'', ``temporary'', etc.
\end{itemize}

Severity is assigned as: \texttt{high} (intensifier present), \texttt{low} (diminisher present), or \texttt{medium} (neither). Risk intensity is computed as a weighted sum: $\text{intensity} = 3 \times n_{\text{high}} + 2 \times n_{\text{medium}} + 1 \times n_{\text{low}}$.

\subsubsection{Context Snippets}

Each detection records a 100-character context snippet (50 characters before and after the keyword) for downstream explainability.

\subsection{Layer 2: Semantic Risk Detection}

The semantic layer addresses the fundamental limitation of keyword approaches: risks expressed using non-standard language are missed entirely. The semantic detector uses sentence-transformer embeddings (\texttt{all-MiniLM-L6-v2}, 384 dimensions) and a differential scoring mechanism.

\subsubsection{Risk and Safe Anchors}

For each of the 10 risk categories, a set of \textbf{risk anchor sentences} (6 per category) and \textbf{safe anchor sentences} (4 per category) are defined, totalling 100 anchor sentences. Risk anchors represent prototypical risk-bearing statements (e.g., ``Customer demand is declining and we are losing subscribers''), while safe anchors represent neutral or positive mentions of the same topic (e.g., ``Demand is strong and growing consistently'').

Risk and safe centroids are computed as the mean embedding of each category's anchors.

\subsubsection{Differential Scoring}

For each text chunk, the system computes:
\begin{equation}
    \text{risk\_similarity} = \cos(\mathbf{e}_{\text{chunk}}, \mathbf{c}_{\text{risk}})
    \label{eq:risk-sim}
\end{equation}
\begin{equation}
    \text{safe\_similarity} = \cos(\mathbf{e}_{\text{chunk}}, \mathbf{c}_{\text{safe}})
    \label{eq:safe-sim}
\end{equation}
\begin{equation}
    \text{differential} = \text{risk\_similarity} - \text{safe\_similarity}
    \label{eq:differential}
\end{equation}

where $\mathbf{e}_{\text{chunk}}$ is the chunk embedding, $\mathbf{c}_{\text{risk}}$ and $\mathbf{c}_{\text{safe}}$ are risk and safe centroids respectively. A chunk is flagged as a risk if both $\text{risk\_similarity} > \text{min\_sim}$ and $\text{differential} > \text{min\_diff}$ for a given category.

\subsubsection{Per-Category Thresholds}

Thresholds were tuned empirically across 49 companies with manual validation. The final thresholds (v5) are:

\begin{table}[H]
\centering
\caption{Semantic risk detection thresholds per category.}
\label{tab:semantic-thresholds}
\small
\begin{tabular}{@{}lcc@{}}
\toprule
\textbf{Category} & \textbf{min\_sim} & \textbf{min\_diff} \\
\midrule
FX / Macro Risk & 0.36 & 0.02 \\
Geopolitical Risk & 0.35 & 0.05 \\
Margin Risk & 0.40 & 0.05 \\
Liquidity Risk & 0.40 & 0.05 \\
Regulatory Risk & 0.33 & 0.07 \\
Tech Execution Risk & 0.38 & 0.07 \\
Supply Chain Risk & 0.37 & 0.08 \\
Demand Risk & 0.40 & 0.08 \\
Competition Risk & 0.40 & 0.10 \\
Labour Risk & 0.38 & 0.10 \\
\bottomrule
\end{tabular}
\end{table}

Lower \texttt{min\_diff} thresholds for FX/macro and geopolitical risks reflect the observation that these topics are frequently discussed in neutral contexts (e.g., ``we hedge our FX exposure''), requiring a more lenient differential to capture genuine risks.

\subsubsection{Explainability via Matched Anchors}

For each semantic detection, the system identifies the closest matching individual anchor sentence by computing cosine similarity between the chunk embedding and each of the category's risk anchor embeddings. The matched anchor is returned alongside the detection, providing a human-readable explanation of \emph{why} the system flagged the chunk.

\subsection{Two-Layer Merging}

When both layers detect the same risk category in the same chunk, the semantic detection is recorded as a \textbf{semantic confirmation} rather than a duplicate. The system tracks three detection types:
\begin{itemize}[noitemsep]
    \item \textbf{Keyword}: Detected by keyword matching only
    \item \textbf{Semantic New}: Detected by semantic similarity in a category with no keyword match
    \item \textbf{Semantic Confirmation}: Semantic detection in a category already found by keywords
\end{itemize}

This distinction enables the ``RAG Value-Add'' metric: the percentage of risk categories that were \emph{only} detected by semantic analysis, measuring the retrieval-augmented detection's contribution beyond keyword matching.


% ══════════════════════════════════════════════════════════════
\section{NLP Feature Engineering}
\label{sec:features}
% ══════════════════════════════════════════════════════════════

A total of 39 NLP features are engineered per earnings event for downstream ML modelling. These are grouped into five categories.

\subsection{Sentiment Features (14 features)}

Direct outputs from the three sentiment models, aggregated to the call level:
\begin{itemize}[noitemsep]
    \item LM: \texttt{mean\_lm\_net\_score}, \texttt{mean\_lm\_positive\_score}, \texttt{mean\_lm\_negative\_score}, \texttt{mean\_lm\_uncertainty\_score}
    \item VADER: \texttt{mean\_vader\_compound}
    \item FinBERT: \texttt{finbert\_positive}, \texttt{finbert\_negative}, \texttt{finbert\_neutral}
    \item Role-based: \texttt{mgmt\_sentiment}, \texttt{analyst\_sentiment}, \texttt{mgmt\_analyst\_divergence}
    \item Section-based: \texttt{prepared\_sentiment}, \texttt{qa\_sentiment}, \texttt{sentiment\_divergence}
\end{itemize}

\subsection{Readability Features (5 features)}

\begin{itemize}[noitemsep]
    \item \textbf{Gunning Fog Index}: $0.4 \times (\text{avg\_sent\_len} + 100 \times \text{complex\_word\_pct})$, where complex words have $\geq 3$ syllables.
    \item \textbf{Flesch--Kincaid Grade Level}: $0.39 \times \text{avg\_sent\_len} + 11.8 \times \text{avg\_syllables} - 15.59$.
    \item \textbf{Average sentence length}, \textbf{average word length}, and \textbf{vocabulary richness} (type-token ratio).
\end{itemize}

\subsection{Forward-Looking Statement Features (2 features)}

\begin{itemize}[noitemsep]
    \item \texttt{fls\_ratio}: Proportion of sentences containing forward-looking keywords (31 terms: ``will'', ``expect'', ``guidance'', ``outlook'', ``forecast'', ``anticipate'', etc.).
    \item \texttt{fls\_count}: Total keyword matches.
\end{itemize}

\subsection{Specificity Features (4 features)}

Quantitative specificity is measured by counting numeric expressions per sentence:
\begin{equation}
    \text{specificity\_score} = \frac{n_{\text{numbers}} + 2 \times n_{\text{dollars}} + 2 \times n_{\text{percents}} + 2 \times n_{\text{bps}}}{n_{\text{sentences}}}
\end{equation}

Dollar amounts and percentages are double-weighted as they represent more precise quantitative commitments.

\subsection{Q\&A Dynamics Features (10 features)}

These features capture behavioural signals from the Q\&A section, motivated by \citeauthor{price2012earnings}'s finding that Q\&A contains more price-relevant information:

\begin{itemize}[noitemsep]
    \item \textbf{Deflection rate}: Fraction of management responses containing evasive phrases (29 patterns: ``as I mentioned'', ``let me come back to'', ``too early to tell'', ``I don't want to speculate'').
    \item \textbf{Filler rate}: Filler word frequency (``um'', ``uh'', ``you know'', ``sort of'', ``basically'').
    \item \textbf{Passive voice rate}: Passive constructions per sentence, indicating diffused accountability.
    \item \textbf{Pronoun shift}: Change in first-person pronoun ratio from prepared remarks to Q\&A. A shift toward third-person (``the company'', ``they'') may signal psychological distancing.
    \item \textbf{Specificity drop}: Decrease in numeric density from prepared remarks to Q\&A, indicating reduced precision under questioning.
    \item \textbf{Answer-question ratio}: Average management answer length relative to analyst question length.
    \item \textbf{Management word ratio}: Proportion of Q\&A words spoken by management.
\end{itemize}

\subsection{Engineered Interaction Features (4 features)}

Research-backed composite features:
\begin{itemize}[noitemsep]
    \item \textbf{Q\&A-Weighted Sentiment} \citep{price2012earnings}: $0.65 \times \text{qa\_sentiment} + 0.35 \times \text{prepared\_sentiment}$
    \item \textbf{Negative Emphasis} \citep{loughran2011liability}: $1.5 \times \text{lm\_negative\_score}$
    \item \textbf{Uncertainty-Adjusted Sentiment}: $\text{net\_score} - 0.5 \times \text{uncertainty\_score}$
    \item \textbf{Risk--Sentiment Interaction}: $\text{net\_score} \times \text{total\_risk\_count}$
\end{itemize}


% ══════════════════════════════════════════════════════════════
\section{RAG System}
\label{sec:rag}
% ══════════════════════════════════════════════════════════════

The Retrieval-Augmented Generation (RAG) subsystem enables natural language querying of earnings call content. It follows a four-stage pipeline: indexing, retrieval, re-ranking, and generation.

\subsection{Indexing}

Text chunks are embedded using \texttt{all-MiniLM-L6-v2} (384-dimensional sentence embeddings) in batches of 64. Embeddings are L2-normalised so that inner product equals cosine similarity, then stored in a FAISS \texttt{IndexFlatIP} index containing 30,235 vectors.

A parallel BM25 index (\texttt{BM25Okapi} from \texttt{rank\_bm25}) is built over the same chunks using lowercase whitespace-tokenised text.

\subsection{Hybrid Retrieval with Reciprocal Rank Fusion}

Given a user query, the system performs both dense (FAISS) and sparse (BM25) search, each retrieving $3 \times k$ candidates (where $k$ is the desired number of results, default 10). Results are fused using Reciprocal Rank Fusion (RRF) \citep{cormack2009reciprocal}:

\begin{equation}
    \text{RRF}(d) = \sum_{r \in \{dense, sparse\}} \frac{w_r}{k_{\text{RRF}} + \text{rank}_r(d)}
    \label{eq:rrf}
\end{equation}

where $k_{\text{RRF}} = 60$ (standard constant), $w_{\text{dense}} = 0.6$, and $w_{\text{sparse}} = 0.4$. The asymmetric weighting reflects the generally superior performance of dense retrieval for semantic queries, while preserving BM25's advantage for exact keyword matching.

\subsection{Cross-Encoder Re-Ranking}

The top $3k$ RRF candidates are re-ranked by a cross-encoder model (\texttt{cross-encoder/ms-marco-MiniLM-L-6-v2}), which jointly encodes the query and each candidate passage to produce a relevance score. This two-stage approach balances efficiency (bi-encoder for candidate retrieval) with precision (cross-encoder for final ranking) \citep{nogueira2019passage}.

\subsection{Answer Generation}

Retrieved chunks are formatted as numbered context passages and fed to Llama~3.2 (3B parameters) running locally via Ollama. The system prompt enforces:
\begin{itemize}[noitemsep]
    \item Use only provided context; do not rely on prior knowledge.
    \item Cite sources using [1], [2], etc.
    \item Distinguish management vs.\ analyst statements.
    \item Note speaker roles when relevant.
\end{itemize}

Generation parameters: \texttt{temperature}=0.1, \texttt{max\_tokens}=1500. A context window limit of 4,096 tokens is enforced by truncating source passages.

\subsection{Query Expansion}

For vague or short queries, the system applies pattern-based query expansion using 6 expansion templates:
\begin{itemize}[noitemsep]
    \item ``how did X do'' $\to$ sub-queries for revenue, earnings, guidance
    \item ``risk'' $\to$ challenge, headwind, pressure, decline
    \item ``strategy'' $\to$ invest, growth, initiative, opportunity
\end{itemize}

Multi-query retrieval deduplicates results by \texttt{chunk\_id} before re-ranking.

\subsection{Citation Tracking}

Citations in the form \texttt{[n]} are extracted from the generated answer via regex and mapped back to the corresponding source chunks, preserving full metadata (ticker, quarter, speaker, role, section) for display in the dashboard.


% ══════════════════════════════════════════════════════════════
\section{Financial Analysis}
\label{sec:finance}
% ══════════════════════════════════════════════════════════════

\subsection{Market Data and Abnormal Returns}

Stock prices are sourced from Yahoo Finance via \texttt{yfinance} with auto-adjustment for splits and dividends. Abnormal returns are computed relative to the S\&P~500 benchmark (SPY):

\begin{equation}
    AR_{i,t} = R_{i,t} - R_{\text{SPY},t}
    \label{eq:abnormal-return}
\end{equation}

Cumulative Abnormal Returns (CAR) are computed over windows of 1, 3, and 5 trading days following the earnings announcement. Abnormal volume is computed as the ratio of earnings-day volume to the trailing 20-day average.

Market data is cached as Parquet files (\texttt{data/market/cache/\{ticker\}\_\{start\}\_\{end\}.parquet}) to minimise API calls, with 0.5-second rate limiting per ticker.

\subsection{Event Study}

A standard event study tests whether post-earnings abnormal returns are statistically different from zero using one-sample $t$-tests. Results on 13,478 events:

\begin{table}[H]
\centering
\caption{Event study results: Cumulative Abnormal Returns.}
\label{tab:event-study}
\begin{tabular}{@{}lcccc@{}}
\toprule
\textbf{Window} & \textbf{Mean CAR} & \textbf{Std Dev} & \textbf{$t$-statistic} & \textbf{$p$-value} \\
\midrule
1-day & 0.104\% & 6.08\% & 1.989 & 0.047* \\
3-day & 0.298\% & 7.64\% & 4.524 & $<$0.001*** \\
5-day & 0.257\% & 8.44\% & 3.529 & $<$0.001*** \\
\bottomrule
\end{tabular}
\end{table}

The positive mean CARs confirm that earnings announcements carry net positive information surprises on average.

\subsection{OLS Regression}

An OLS regression of 1-day abnormal returns on NLP features yields:

\begin{table}[H]
\centering
\caption{OLS regression results (dependent variable: 1-day abnormal return).}
\label{tab:ols}
\small
\begin{tabular}{@{}lccc@{}}
\toprule
\textbf{Feature} & \textbf{Coefficient} & \textbf{$t$-stat} & \textbf{$p$-value} \\
\midrule
Analyst sentiment & 0.459 & 3.66 & $<$0.001*** \\
Sentiment momentum (LM) & 0.452 & 3.86 & $<$0.001*** \\
Sentiment divergence & 0.171 & 3.14 & 0.002** \\
Mgmt--analyst divergence & $-$0.151 & $-$2.65 & 0.008** \\
Sentiment momentum (VADER) & 0.030 & 3.18 & 0.002** \\
\bottomrule
\end{tabular}
\end{table}

The model achieves $R^2 = 0.017$ ($F = 7.69$, $p < 0.001$). While the $R^2$ is low in absolute terms, it is consistent with the finance literature where even 1--2\% explanatory power from text features is considered economically meaningful \citep{chen2018sentimental}. Notably, \textbf{analyst sentiment} and \textbf{sentiment momentum} are the strongest predictors, consistent with \citeauthor{price2012earnings}'s findings on Q\&A informativeness.

\subsection{Portfolio Sorts}

Earnings events are sorted into terciles by sentiment score, and average post-earnings abnormal returns are computed per bucket:

\begin{table}[H]
\centering
\caption{Portfolio sort results by LM net sentiment score.}
\label{tab:portfolio-sort}
\begin{tabular}{@{}lccc@{}}
\toprule
\textbf{Bucket} & \textbf{$n$} & \textbf{Mean Return} & \textbf{Std Dev} \\
\midrule
Low sentiment & 4,493 & $-$0.252\% & 5.90\% \\
Mid sentiment & 4,492 & +0.133\% & 6.17\% \\
High sentiment & 4,493 & +0.432\% & 6.16\% \\
\midrule
\textbf{Long--Short Spread} & & \textbf{+0.684\%} & \\
\bottomrule
\end{tabular}
\end{table}

The monotonic increase from low to high sentiment terciles, and the positive long-short spread (72 basis points using VADER compound), demonstrates that text-derived sentiment carries economically meaningful information for systematic trading strategies.

\subsection{ML Models for Direction Prediction}

Three classifiers are trained to predict the binary direction of post-earnings abnormal returns ($AR > 0$):

\begin{table}[H]
\centering
\caption{Direction prediction: 5-fold stratified cross-validation accuracy.}
\label{tab:ml-accuracy}
\begin{tabular}{@{}lcc@{}}
\toprule
\textbf{Model} & \textbf{Accuracy} & \textbf{Std Dev} \\
\midrule
Logistic Regression & 53.1\% & 0.65\% \\
Random Forest (200 trees, depth 5) & 52.5\% & 0.71\% \\
Gradient Boosting (200 trees, depth 3) & 51.9\% & 0.67\% \\
\midrule
Baseline (majority class) & 50.9\% & --- \\
\bottomrule
\end{tabular}
\end{table}

Additionally, four regression models are evaluated using walk-forward time-series cross-validation (5-fold \texttt{TimeSeriesSplit}):

\begin{table}[H]
\centering
\caption{Return regression: walk-forward cross-validation $R^2$.}
\label{tab:ml-r2}
\begin{tabular}{@{}lcc@{}}
\toprule
\textbf{Model} & \textbf{Mean $R^2$} & \textbf{Std Dev} \\
\midrule
Ridge ($\alpha = 1.0$) & 0.0103 & 0.0052 \\
Random Forest (100 trees, depth 3) & 0.0081 & 0.0038 \\
Lasso ($\alpha = 0.01$) & $-$0.0005 & 0.0006 \\
Gradient Boosting (100 trees, depth 2) & $-$0.0001 & 0.0073 \\
\bottomrule
\end{tabular}
\end{table}

The modest classification accuracy (2 percentage points above baseline) and near-zero regression $R^2$ are \textbf{consistent with the academic literature}. \citet{chen2018sentimental} report that NLP features alone have limited out-of-sample predictive power for individual stock returns, as prices rapidly incorporate textual information. The economic value lies in systematic portfolio strategies (see \Cref{tab:portfolio-sort}) rather than individual-event prediction.


% ══════════════════════════════════════════════════════════════
\section{Evaluation}
\label{sec:evaluation}
% ══════════════════════════════════════════════════════════════

\subsection{Retrieval Evaluation}

Retrieval quality is evaluated on 12 hand-crafted test queries spanning 10 companies, covering diverse information needs (revenue guidance, risk factors, competitive strategy, analyst concerns, CFO commentary). Each query specifies expected tickers, keywords, sections, and roles.

\begin{table}[H]
\centering
\caption{Retrieval evaluation metrics.}
\label{tab:retrieval-eval}
\begin{tabular}{@{}lcc@{}}
\toprule
\textbf{Metric} & \textbf{Score} & \textbf{Std Dev} \\
\midrule
Recall@3 & 86.9\% & 18.9\% \\
Recall@5 & 86.9\% & 18.9\% \\
Recall@10 & 95.8\% & 13.8\% \\
MRR (Mean Reciprocal Rank) & 100.0\% & 0.0\% \\
Ticker Accuracy & 91.7\% & 15.2\% \\
Keyword Hit Rate & 95.8\% & 13.8\% \\
\bottomrule
\end{tabular}
\end{table}

The perfect MRR (1.0) indicates that the most relevant chunk is always ranked first after hybrid retrieval and cross-encoder re-ranking.

\subsection{Answer Quality Evaluation}

Answer quality is assessed on 6 queries with automated metrics:

\begin{table}[H]
\centering
\caption{Answer quality metrics.}
\label{tab:answer-eval}
\begin{tabular}{@{}lc@{}}
\toprule
\textbf{Metric} & \textbf{Score} \\
\midrule
Success Rate & 100\% (6/6) \\
Average Answer Length & 85.8 words \\
Citation Rate & 83.3\% \\
Keyword Coverage & 83.3\% \\
Hedge Rate & 0.0\% \\
Query Expansion Rate & 33.3\% \\
Average Sources per Query & 5.0 \\
Average Unique Speakers & 1.67 \\
\bottomrule
\end{tabular}
\end{table}

\subsection{Semantic Risk Detection Validation}

Semantic risk thresholds underwent five iterations of tuning (v1--v5). Initial thresholds were too conservative, resulting in zero semantic detections on real transcripts. The tuning process involved:
\begin{enumerate}[noitemsep]
    \item Running diagnostic analyses on test transcripts to compute per-category similarity and differential scores for each chunk.
    \item Identifying chunks where scores fell just below thresholds, indicating the thresholds were too strict.
    \item Progressively relaxing \texttt{min\_diff} from 0.07--0.16 (v1) to 0.02--0.10 (v5).
    \item Validating that relaxed thresholds did not introduce false positives on clearly non-risk text.
\end{enumerate}


% ══════════════════════════════════════════════════════════════
\section{Interactive Dashboard}
\label{sec:dashboard}
% ══════════════════════════════════════════════════════════════

The system is deployed as a multi-page Streamlit application with five specialised pages:

\begin{enumerate}[noitemsep]
    \item \textbf{Home}: PDF upload with auto-detection of ticker and quarter, multi-file support, and analysis orchestration.
    \item \textbf{Sentiment} (\texttt{01\_Sentiment.py}): Three-model sentiment comparison, management--analyst gap visualisation, credibility waterfall (prepared $\to$ Q\&A), quarter-over-quarter momentum charts, and a model glossary.
    \item \textbf{Risk} (\texttt{02\_Risk.py}): Key findings summary box, detection method breakdown (keyword vs.\ semantic new vs.\ semantic confirmation), management vs.\ analyst risk exposure, and top risk excerpts with method badges and matched anchor annotations.
    \item \textbf{Benchmark} (\texttt{03\_Benchmark.py}): ML direction prediction with plain-English interpretation, per-quarter NLP feature tabs, portfolio sort visualisation, and sentiment momentum tracking.
    \item \textbf{Transcript} (\texttt{04\_Transcript.py}): Parsed transcript viewer with speaker annotations and section breakdown.
    \item \textbf{Q\&A} (\texttt{05\_QA.py}): RAG-powered question answering with suggested questions, source citations, multi-quarter search, and speaker/section filters.
\end{enumerate}

The dashboard is styled with a custom dark theme (\texttt{style.css}) featuring gradient accents, hover animations, and responsive typography. All pages include plain-English key findings boxes to support non-technical users.

Multi-quarter mode enables uploading multiple transcripts for the same company. The system auto-detects tickers and quarters, validates company consistency, flags temporal gaps, and enables cross-quarter analysis on all pages.


% ══════════════════════════════════════════════════════════════
\section{Design Decisions and Rationale}
\label{sec:decisions}
% ══════════════════════════════════════════════════════════════

\subsection{Local-First Architecture}

All models run on-device (FinBERT, sentence-transformers, Llama~3.2 via Ollama) with zero API dependencies. This ensures:
\begin{itemize}[noitemsep]
    \item \textbf{Data privacy}: Transcripts---potentially containing material non-public information---never leave the machine.
    \item \textbf{Cost}: No per-token API charges, enabling analysis of thousands of transcripts.
    \item \textbf{Reproducibility}: No dependence on external model versions or rate limits.
\end{itemize}

\subsection{Three-Model Sentiment Ensemble}

Using three sentiment models rather than one addresses the complementary strengths and weaknesses of each approach:
\begin{itemize}[noitemsep]
    \item LM captures \emph{domain-specific} terminology missed by general models.
    \item VADER captures \emph{intensity} and \emph{emphasis} through punctuation and modifier rules.
    \item FinBERT captures \emph{contextual semantics} (negation, sarcasm, implicit sentiment).
\end{itemize}

\subsection{Two-Layer Risk Detection}

The keyword layer provides high-precision, interpretable detections, while the semantic layer provides recall for non-standard risk language. The merged output with explicit detection-method labels gives users transparency into \emph{how} each risk was identified.

\subsection{Differential Scoring for Semantic Risks}

Naive cosine similarity against risk anchors alone produces high false-positive rates because many financial topics are discussed in both risky and safe contexts. The differential score (risk similarity minus safe similarity) effectively filters neutral mentions, improving precision without sacrificing recall.

\subsection{Hybrid Retrieval}

Dense retrieval (FAISS) excels at semantic matching but struggles with exact terms (ticker symbols, specific dollar amounts). BM25 complements this with precise keyword matching. RRF fusion provides a robust, parameter-light combination method \citep{cormack2009reciprocal}.

\subsection{Honest ML Reporting}

The ML direction accuracy of 53\% is reported transparently alongside the baseline of 50.9\%. Rather than overfitting to inflate accuracy, the system emphasises economically meaningful signals (portfolio sort spreads) and acknowledges the inherent difficulty of predicting individual stock returns from text alone---consistent with the academic consensus \citep{chen2018sentimental}.


% ══════════════════════════════════════════════════════════════
\section{Limitations and Future Work}
\label{sec:limitations}
% ══════════════════════════════════════════════════════════════

\subsection{Limitations}

\begin{itemize}[noitemsep]
    \item \textbf{Speaker detection accuracy}: Regex-based parsing is brittle across different transcript formatting conventions. Some transcripts from less common providers may have misattributed speakers.
    \item \textbf{Semantic risk thresholds}: Thresholds were tuned on S\&P~500 transcripts and may not generalise to small-cap companies or non-English transcripts.
    \item \textbf{ML predictive power}: Direction accuracy is only marginally above baseline. NLP features alone are insufficient for profitable individual-event trading.
    \item \textbf{LLM hallucination}: Despite system prompt constraints, Llama~3.2 may occasionally generate information not present in the retrieved context.
    \item \textbf{Evaluation set size}: Retrieval evaluation uses only 12 queries, which limits statistical power.
\end{itemize}

\subsection{Future Work}

\begin{itemize}[noitemsep]
    \item \textbf{Larger evaluation sets}: Expand retrieval and answer quality evaluation to 100+ queries with human annotation.
    \item \textbf{Fine-tuned risk detector}: Train a supervised classifier on labelled risk segments rather than relying on anchor-based similarity.
    \item \textbf{Multimodal features}: Incorporate audio features (tone, pace, hesitation) from earnings call recordings.
    \item \textbf{Real-time processing}: Stream transcripts as they are released for immediate analysis.
    \item \textbf{Cross-company analysis}: Enable comparative analysis across peers within the same sector.
\end{itemize}


% ══════════════════════════════════════════════════════════════
\section{Conclusion}
\label{sec:conclusion}
% ══════════════════════════════════════════════════════════════

This project demonstrates that a fully local, end-to-end NLP system can extract meaningful financial intelligence from earnings call transcripts. The three-model sentiment pipeline captures complementary dimensions of language, while the two-layer risk detection framework provides both precision (keyword matching) and recall (semantic embeddings) with full explainability. The hybrid RAG system achieves near-perfect retrieval metrics (Recall@3 = 86.9\%, MRR = 100\%) and generates grounded, cited answers. While ML direction prediction is marginal, portfolio sorts demonstrate a statistically significant and economically meaningful long-short spread of 72 basis points per trade across 13,478 events. The interactive Streamlit dashboard makes these capabilities accessible to both technical and non-technical users, all running locally with zero API costs and complete data privacy.


% ══════════════════════════════════════════════════════════════
\bibliographystyle{apalike}
\begin{thebibliography}{99}

\bibitem[Araci, 2019]{araci2019finbert}
Araci, D. (2019).
\newblock FinBERT: Financial Sentiment Analysis with Pre-Trained Language Models.
\newblock \emph{arXiv preprint arXiv:1908.10063}.

\bibitem[Chen et~al., 2018]{chen2018sentimental}
Chen, H., De, P., Hu, Y., \& Hwang, B.-H. (2018).
\newblock Wisdom of Crowds: The Value of Stock Opinions Transmitted Through Social Media.
\newblock \emph{The Review of Financial Studies}, 27(5), 1367--1403.

\bibitem[Cormack et~al., 2009]{cormack2009reciprocal}
Cormack, G.~V., Clarke, C.~L., \& Buettcher, S. (2009).
\newblock Reciprocal Rank Fusion Outperforms Condorcet and Individual Rank Learning Methods.
\newblock \emph{Proceedings of the 32nd International ACM SIGIR Conference}, 758--759.

\bibitem[Hutto \& Gilbert, 2014]{hutto2014vader}
Hutto, C.~J. \& Gilbert, E. (2014).
\newblock VADER: A Parsimonious Rule-Based Model for Sentiment Analysis of Social Media Text.
\newblock \emph{Proceedings of the 8th International AAAI Conference on Weblogs and Social Media}.

\bibitem[Loughran \& McDonald, 2011]{loughran2011liability}
Loughran, T. \& McDonald, B. (2011).
\newblock When Is a Liability Not a Liability? Textual Analysis, Dictionaries, and 10-Ks.
\newblock \emph{The Journal of Finance}, 66(1), 35--65.

\bibitem[Nogueira \& Cho, 2019]{nogueira2019passage}
Nogueira, R. \& Cho, K. (2019).
\newblock Passage Re-Ranking with BERT.
\newblock \emph{arXiv preprint arXiv:1901.04085}.

\bibitem[Price et~al., 2012]{price2012earnings}
Price, S.~M., Doran, J.~S., Peterson, D.~R., \& Bliss, B.~A. (2012).
\newblock Earnings Conference Calls and Stock Returns: The Incremental Informativeness of Textual Tone.
\newblock \emph{Journal of Banking \& Finance}, 36(4), 992--1011.

\bibitem[Robertson et~al., 2009]{robertson2009probabilistic}
Robertson, S. \& Zaragoza, H. (2009).
\newblock The Probabilistic Relevance Framework: BM25 and Beyond.
\newblock \emph{Foundations and Trends in Information Retrieval}, 3(4), 333--389.

\bibitem[Reimers \& Gurevych, 2019]{reimers2019sentence}
Reimers, N. \& Gurevych, I. (2019).
\newblock Sentence-BERT: Sentence Embeddings using Siamese BERT-Networks.
\newblock \emph{Proceedings of EMNLP-IJCNLP 2019}, 3982--3992.

\bibitem[Devlin et~al., 2019]{devlin2019bert}
Devlin, J., Chang, M.-W., Lee, K., \& Toutanova, K. (2019).
\newblock BERT: Pre-Training of Deep Bidirectional Transformers for Language Understanding.
\newblock \emph{Proceedings of NAACL-HLT 2019}, 4171--4186.

\bibitem[Johnson et~al., 2019]{johnson2019faiss}
Johnson, J., Douze, M., \& J\'{e}gou, H. (2019).
\newblock Billion-Scale Similarity Search with GPUs.
\newblock \emph{IEEE Transactions on Big Data}, 7(3), 535--547.

\bibitem[Lewis et~al., 2020]{lewis2020rag}
Lewis, P., Perez, E., Piktus, A., et~al. (2020).
\newblock Retrieval-Augmented Generation for Knowledge-Intensive NLP Tasks.
\newblock \emph{Advances in Neural Information Processing Systems}, 33, 9459--9474.

\end{thebibliography}

% ══════════════════════════════════════════════════════════════
\appendix
\section{Complete Feature List}
\label{app:features}
% ══════════════════════════════════════════════════════════════

\begin{longtable}{@{}rllp{6.5cm}@{}}
\caption{All 39 NLP features used in ML models.} \label{tab:all-features} \\
\toprule
\textbf{\#} & \textbf{Feature} & \textbf{Category} & \textbf{Description} \\
\midrule
\endfirsthead
\toprule
\textbf{\#} & \textbf{Feature} & \textbf{Category} & \textbf{Description} \\
\midrule
\endhead
1 & mean\_lm\_net\_score & Sentiment & LM positive minus negative proportion \\
2 & mean\_lm\_positive\_score & Sentiment & LM positive word proportion \\
3 & mean\_lm\_negative\_score & Sentiment & LM negative word proportion \\
4 & mean\_lm\_uncertainty\_score & Sentiment & LM uncertainty word proportion \\
5 & mean\_vader\_compound & Sentiment & VADER compound score [$-$1, 1] \\
6 & finbert\_positive & Sentiment & FinBERT positive probability \\
7 & finbert\_negative & Sentiment & FinBERT negative probability \\
8 & finbert\_neutral & Sentiment & FinBERT neutral probability \\
9 & mgmt\_sentiment & Sentiment & Management speakers' avg LM score \\
10 & analyst\_sentiment & Sentiment & Analyst speakers' avg LM score \\
11 & prepared\_sentiment & Sentiment & Prepared remarks avg LM score \\
12 & qa\_sentiment & Sentiment & Q\&A section avg LM score \\
13 & mgmt\_analyst\_divergence & Sentiment & mgmt minus analyst sentiment \\
14 & sentiment\_divergence & Sentiment & prepared minus Q\&A sentiment \\
15 & fog\_index & Readability & Gunning Fog readability index \\
16 & flesch\_kincaid & Readability & Flesch--Kincaid grade level \\
17 & avg\_sentence\_length & Readability & Mean words per sentence \\
18 & avg\_word\_length & Readability & Mean characters per word \\
19 & vocab\_richness & Readability & Type-token ratio \\
20 & fls\_ratio & FLS & Forward-looking statement proportion \\
21 & fls\_count & FLS & Forward-looking keyword count \\
22 & specificity\_score & Specificity & Weighted numeric density per sentence \\
23 & numbers\_per\_sentence & Specificity & Numeric mentions per sentence \\
24 & dollar\_mentions & Specificity & Dollar amount count \\
25 & percent\_mentions & Specificity & Percentage mention count \\
26 & qa\_mgmt\_word\_ratio & Q\&A & Management word proportion in Q\&A \\
27 & qa\_avg\_answer\_length & Q\&A & Mean management answer length \\
28 & qa\_avg\_question\_length & Q\&A & Mean analyst question length \\
29 & qa\_answer\_question\_ratio & Q\&A & Answer / question length ratio \\
30 & qa\_deflection\_rate & Q\&A & Evasion phrase frequency \\
31 & qa\_filler\_rate & Q\&A & Filler word frequency \\
32 & qa\_passive\_rate & Q\&A & Passive voice construction rate \\
33 & qa\_pronoun\_shift & Q\&A & First-person pronoun ratio change \\
34 & qa\_specificity\_drop & Q\&A & Numeric density decrease in Q\&A \\
35 & qa\_weighted\_sentiment & Engineered & $0.65 \times$ Q\&A + $0.35 \times$ prepared \\
36 & negative\_emphasis & Engineered & $1.5 \times$ LM negative score \\
37 & neg\_pos\_ratio & Engineered & Negative / positive score ratio \\
38 & uncertainty\_adj\_sentiment & Engineered & Net score $-$ $0.5 \times$ uncertainty \\
39 & risk\_sentiment\_interaction & Engineered & Net score $\times$ risk count \\
\bottomrule
\end{longtable}

\section{Risk Taxonomy Keywords}
\label{app:risk-keywords}

\begin{longtable}{@{}lp{11cm}@{}}
\caption{Risk categories with full keyword lists (255 total).} \label{tab:risk-keywords} \\
\toprule
\textbf{Category} & \textbf{Keywords} \\
\midrule
\endfirsthead
\toprule
\textbf{Category} & \textbf{Keywords} \\
\midrule
\endhead
Demand Risk & declining demand, weak demand, order cancellations, pipeline weakness, lower bookings, softening demand, demand uncertainty, customer churn, volume decline, revenue decline, reduced backlog, deteriorating demand, lower orders, demand softness, sales decline, shrinking market, market contraction, consumption decline, booking decline, subscriber loss, traffic decline, footfall decline, fewer transactions, purchase delays, project deferrals, deal slippage, delayed orders, elongated sales cycle \\
\midrule
Margin Risk & margin pressure, cost inflation, pricing pressure, gross margin decline, operating margin pressure, input cost increase, wage inflation, freight cost, margin erosion, margin compression, cost overrun, unfavorable mix, dilutive impact, promotional activity, markdown, cost headwind, shrinkage, underabsorption, negative operating leverage, raw material cost increase, elevated costs, pricing headwind, lower margin, cost pressure, higher costs, profitability decline \\
\midrule
Supply Chain Risk & supply chain disruption, component shortage, inventory buildup, capacity constraints, logistics challenges, lead time, supplier risk, raw material shortage, transportation disruption, warehouse capacity, port congestion, shipping delay, procurement risk, single source, manufacturing bottleneck, supply shortage, inventory excess, distribution issues, fulfillment delays, sourcing challenges, stockout, production delay, logistics bottleneck, backlog increase \\
\midrule
Regulatory Risk & regulatory, litigation, antitrust, GDPR, data privacy, compliance cost, regulatory uncertainty, tax audit, consent decree, class action, patent dispute, investigation, enforcement action, penalty, fine, regulatory change, legislative risk, policy uncertainty, FDA, EPA, SEC inquiry, regulatory scrutiny, legal proceedings, license revocation, regulatory burden, import restriction, export control, government investigation \\
\midrule
Competition Risk & competitive pressure, market share loss, commoditization, disruptive competition, new entrant, price competition, competitive landscape, aggressive competitor, competitive threat, competitive dynamics, market saturation, intensifying competition, loss of customers, switching cost, competitive disadvantage, competitive erosion, competitor pricing, rivalry, market fragmentation, differentiation challenge, brand erosion, competitive intensity, losing ground, share loss \\
\midrule
FX / Macro Risk & foreign exchange headwind, currency impact, strong dollar, recession risk, inflation pressure, interest rate risk, monetary tightening, credit tightening, economic slowdown, GDP decline, consumer weakness, FX headwind, currency volatility, exchange rate impact, macro uncertainty, economic uncertainty, stagflation, yield curve, rate hike, fiscal tightening, purchasing power decline, consumer confidence decline, housing slowdown, credit crunch, macro deterioration, macro headwind, inflationary pressure \\
\midrule
Liquidity Risk & cash burn, negative free cash flow, debt maturity, covenant compliance, credit rating downgrade, refinancing risk, liquidity constraint, working capital pressure, cash flow deterioration, debt covenant, credit facility, capital raise, dilution, funding gap, accounts receivable aging, cash position, debt service, leverage ratio, solvency concern, cash shortfall, borrowing cost, financial flexibility, debt load, debt reduction \\
\midrule
Tech Execution Risk & product delay, cybersecurity breach, technical debt, scalability issues, system outage, integration risk, platform migration, data breach, software defect, legacy system, technology obsolescence, IT failure, security vulnerability, cloud migration risk, implementation delay, development delay, quality issue, product recall, bug, downtime, system failure, service disruption, reliability issue, modernization challenge \\
\midrule
Labour Risk & talent shortage, employee attrition, key person risk, labour dispute, unionization, compensation pressure, hiring freeze, layoffs, restructuring charge, severance cost, headcount reduction, retention challenge, burnout, workforce reduction, recruitment difficulty, skill gap, staffing challenge, labour cost increase, strike, workplace safety, absenteeism, turnover rate, talent retention, succession risk \\
\midrule
Geopolitical Risk & trade war, tariffs, China risk, sanctions, political instability, war, conflict, geopolitical tension, export restrictions, import restrictions, trade dispute, embargo, political risk, sovereign risk, country risk, expropriation, currency control, capital control, political uncertainty, decoupling, supply chain repatriation, national security, foreign policy, territorial dispute, regime change, diplomatic tension \\
\bottomrule
\end{longtable}


\section{Evaluation Query Set}
\label{app:eval-queries}

\begin{longtable}{@{}clp{8cm}@{}}
\caption{Retrieval evaluation queries (12 total).} \label{tab:eval-queries} \\
\toprule
\textbf{\#} & \textbf{Ticker} & \textbf{Query} \\
\midrule
\endfirsthead
1 & AAPL & What did Apple say about iPhone revenue? \\
2 & MSFT & How is Microsoft's cloud business performing? \\
3 & NVDA & What is NVIDIA's data center growth outlook? \\
4 & TSLA & What are Tesla's production and delivery challenges? \\
5 & META & How is Meta's advertising revenue trending? \\
6 & AMZN & What did Amazon report about AWS margins? \\
7 & NFLX & What are Netflix's subscriber growth numbers? \\
8 & CRM & What is Salesforce's revenue guidance? \\
9 & ORCL & How are Oracle's cloud operations performing? \\
10 & MSFT & What competitive threats does Microsoft face? \\
11 & --- & What did the CFO say about operating expenses? \\
12 & --- & What concerns did analysts raise? \\
\bottomrule
\end{longtable}


\end{document}
